% Generated from locale/tr/content/computability/recursive-functions/introduction.tex; do not hand-edit.


\olfileid[tr]{cmp}{rec}{int}
\olsection{Giriş}

Matematiksel bir hesaplanabilirlik kuramı geliştirmek için öncelikle
bir hesaplanabilirlik \emph{modeli} geliştirmek gerekir. Bugün
hesaplanabilirliği bilgisayarların yaptığı türden bir iş olarak
düşünürüz ve bilgisayarlar simgelerle çalışır. Fakat hesaplanabilirlik
kuramlarının gelişiminin başlangıcında hesaplamanın örnek örneği
\emph{sayısal} hesaplamaydı. Matematikçiler her zaman hesaplanabilen
sayı kuramsal fonksiyonlarla, yani $f\colon\Nat^n\to\Nat$
fonksiyonlarıyla ilgilenmiştir. Dolayısıyla hesaplanabilirlik kuramının
başlangıcında bu tür fonksiyonların incelenmesi şaşırtıcı değildir.
Toplama, çarpma ve doğal sayıların üs alma işlemi gibi en tanıdık
hesaplanabilir sayısal fonksiyonlar ilginç bir özelliği paylaşır:
\emph{özyinelemeyle} tanımlanabilirler. Bu nedenle hesaplanabilir
fonksiyonun genel bir tanımını özyinelemeli tanımlara dayandırmayı
denemek oldukça doğaldır. Sayı kuramsal fonksiyonları özyinelemeyle
tanımlamanın birçok olası yolu arasında özellikle basit bir tanım
örüntüsü merkezî hâle gelir: \emph{ilkel özyineleme}.

Hesaplanabilir fonksiyonların yanı sıra hesaplanabilir küme ve
bağıntılarla da ilgilenebiliriz. Verilen bir sayının bir kümenin
!!a{element} ifadesi olup olmadığının yanıtını hesaplayabiliyorsak o
küme hesaplanabilirdir. Bir $\tuple{n_1,\dots,n_k}$ ikilisinin bir
bağıntının !!a{element} ifadesi olup olmadığını hesaplayabiliyorsak
bağıntı hesaplanabilirdir. Bir küme veya bağıntının \emph{karakteristik
fonksiyonunu} ele alarak hesaplanabilir küme ve bağıntıların
incelenmesini hesaplanabilir fonksiyonların incelenmesine dâhil
edebiliriz. Böylece ilkel özyinelemeli bağıntıları da tanımlayabiliriz;
örneğin ``$n$, $m$'yi tam böler'' bağıntısı ilkel özyinelemeli bir
bağıntıdır.

Yalnızca ilkel özyinelemeyle tanımlanabilen ilkel özyinelemeli
fonksiyonlar, hesaplanabilir sayı kuramsal fonksiyonların tümü değildir.
İlkel özyinelemenin birçok genellemesi ele alınmıştır; fakat fonksiyon
hesaplamanın en güçlü ve en yaygın kabul gören ek yolu sınırsız
aramadır. Bu, \emph{kısmi özyinelemeli fonksiyonlar} tanımına ve onunla
ilişkili \emph{genel özyinelemeli fonksiyonlar} tanımına götürür. Genel
özyinelemeli fonksiyonlar hesaplanabilir ve tümeldir; tanım, tümel olan
kısmi özyinelemeli fonksiyonları tam olarak niteler. Özyinelemeli
fonksiyonlar diğer her hesaplama modelini (Turing makineleri, lambda
hesabı vb.) benzetebilir ve bu yüzden kabul edilen birçok hesaplama
modelinden birini temsil eder.

